%%%%%%%%%%%%%%%%%%%%%%%%%%%%%%%%%%%%%%%%%%%%%%%%%%%%%%%%%%%%%%%%%%%%%%
% Lớp 11 | Xác suất | Biến cố hợp và quy tắc cộng xác suất
% Dạng 1D9?2-4: Quy tắc cộng cho hai biến cố bất kì
% Số câu: 34  (sinh tự động bởi tach_cau.py; từ nay thầy sửa trực tiếp file này)
%%%%%%%%%%%%%%%%%%%%%%%%%%%%%%%%%%%%%%%%%%%%%%%%%%%%%%%%%%%%%%%%%%%%%%

%%%%%%%%%%%%%%%% TRẮC NGHIỆM 4 PHƯƠNG ÁN %%%%%%%%%%%%%%%%

% nguồn: Ok_TracNghiemLop11_Lan1
\begin{ex}%[1D9N2-4]%[ma:5f84f50a]
	Gieo $2$ con xúc xắc cân đối và đồng chất. Xác suất của biến cố \lq\lq Tổng số chấm xuất hiện trên  hai con xúc xắc chia hết cho $5$ \rq\rq\,là
	 \choice
		{$\dfrac{5}{36}$} 
		 {$\dfrac{1}{6}$}
		 {\True $\dfrac{7}{36}$}
		 {$\dfrac{2}{9}$}
	 \loigiai{
		 Kết quả thuận lợi cho biến cố \lq\lq Tổng số chấm xuất hiện trên hai con xúc xắc là $5$ \rq\rq\,là $4$.\\
		Kết  quả thuận lợi cho biến cố \lq\lq Tổng số chấm xuất hiện trên hai con xúc xắc là $10$ \rq\rq\,là $3$.\\
		Kết quả thuận lợi cho biến cố \lq\lq Tổng số chấm xuất hiện trên hai con xúc xắc chia hết cho $5$ \rq\rq\,là  $3+4=7$.\\
	Xác suất của biến cố \lq\lq Tổng số chấm xuất hiện trên hai con xúc xắc chia hết cho $5$  \rq\rq\,là $\dfrac{7}{36}$.}
\end{ex}

% nguồn: Ok_TracNghiemLop11_Lan1
\begin{ex}%[1D9N2-4]%[ma:81433a5a]
	Cho $A$ và $B$ là hai biến cố xung khắc.  Biết $\mathrm{P}(A)=\dfrac{1}{3}$, $\mathrm{P}(B)=\dfrac{1}{4}$. Tính $\mathrm{P}(A\cup B)$.
	\choice 
		{\True $\dfrac{7}{12}$} 
		{$\dfrac{1}{12}$} 
		 {$\dfrac{1}{7}$}
		 {$\dfrac{1}{2}$}
	\loigiai{ 
		Do $A$ và $B$ là  hai biến cố xung khắc nên
	$$\mathrm{P}(A\cup  B)=\mathrm{P}(A)+\mathrm{P}(B)=\dfrac{1}{3}+\dfrac{1}{4}=\dfrac{7}{12}.$$}
\end{ex}

% nguồn: Ok_TracNghiemLop11_Lan1
\begin{ex}%[1D9N2-4]%[ma:6b63104d]
	Cho hai biến cố $A$ và $B$ là hai  biến cố xung khắc. Biết $\mathrm{P}(B)=\dfrac{2}{5}$, $\mathrm{P}(A\cup B)=\dfrac{1}{2}$. Tính $\mathrm{P}(A)$
	 \choice
		 {$\dfrac{9}{10}$}
		{\True $\dfrac{1}{10}$} 
		 {$\dfrac{1}{5}$}
		 {$\dfrac{3}{4}$}
	\loigiai{ 
		Vì $A$  và $B$ là các biến cố xung khắc nên $\mathrm{P}(A\cup B)=\mathrm{P}(A)+\mathrm{P}(B)$.\\
	Suy  ra $\mathrm{P}(A)=\mathrm{P}(A\cup B)-\mathrm{P}(B)=\dfrac{1}{2}-\dfrac{2}{5}=\dfrac{1}{10}$.}
\end{ex}

% nguồn: Ok_TracNghiemLop11_Lan1
\begin{ex}%[1D9N2-4]%[ma:e9472669]
	Cho hai biến cố $A$ và $B$ là hai biến cố xung khắc. Biết $\mathrm{P}(A)=\dfrac{1}{4}$, $\mathrm{P}(A\cup B)=\dfrac{1}{2}$.  Tính $\mathrm{P}(B)$
	 \choice
		{$\dfrac{1}{8}$} 
		{\True  $\dfrac{1}{4}$}
		{$\dfrac{1}{3}$} 
		 {$\dfrac{3}{4}$}
	 \loigiai{
		Vì $A$  và $B$ là các biến cố xung khắc nên $\mathrm{P}(A\cup B)=\mathrm{P}(A)+\mathrm{P}(B)$.\\
	Suy  ra $\mathrm{P}(B)=\mathrm{P}(A\cup B)-\mathrm{P}(A)=\dfrac{1}{2}-\dfrac{1}{4}=\dfrac{1}{4}$.}
\end{ex}

% nguồn: Ok_TracNghiemLop11_Lan1
\begin{ex}%[1D9N2-4]%[ma:ec94cabf]
	Cho $A$, $B$ là hai  biến cố xung khắc. Biết $\mathrm{P}(A)=\dfrac{1}{5}$, $\mathrm{P}(A\cup B)=\dfrac{1}{3}$. Tính $\mathrm{P}(B)$.
	 \choice
		 {$\dfrac{3}{5}$}
		{$\dfrac{8}{15}$} 
		{\True $\dfrac{2}{15}$} 
		 {$\dfrac{1}{15}$}
	 \loigiai{
		$A$, $B$ là  hai biến cố xung khắc.\\
	Suy  ra $\mathrm{P}(A\cup B)=\mathrm{P}(A)+\mathrm{P}(B)\Rightarrow \mathrm{P}(B)=\mathrm{P}(A\cup B)-\mathrm{P}(A)=\dfrac{1}{3}-\dfrac{1}{5}=\dfrac{2}{15}$.}
\end{ex}

% nguồn: Ok_TracNghiemLop11_Lan1
\begin{ex}%[1D9H2-4]%[ma:1edafa39]
	Lớp có $44$ học sinh gồm $24$ học sinh giỏi và $20$ học sinh khá. Thầy giáo chủ nhiệm cần chọn ngẫu nhiên trong lớp một  nhóm gồm $3$ học sinh để kiểm tra kiến thức cũ. Tính xác suất trong $3$ bạn thầy chọn số học sinh giỏi nhiều hơn học sinh khá.
	\choice 
		 {$\dfrac{1}{14}$}
		{\True  $\dfrac{1886}{3311}$}
		{$\dfrac{1380}{3311}$} 
		 {$\dfrac{46}{301}$}
	 \loigiai{
		Ta  có $n(\Omega)=\mathrm{C}_{44}^3$.\\
		 Xét biến cố $A$ là biến cố cần tìm.\\
		Xét  các trường hợp sau
		\begin{itemize}
			\item \textit{Trường hợp  1:} $3$ học sinh giỏi, $0$ học sinh khá có $\mathrm{C}_{24}^3\cdot \mathrm{C}_{20}^0$ cách.
			\item \textit{Trường hợp 2:} $2$ học  sinh giỏi, $1$ học sinh khá có $\mathrm{C}_{24}^2\cdot \mathrm{C}_{20}^1$ cách.
		\end{itemize}
	Suy ra $n(A)=7544$  hay $\mathrm{P}(A)=\dfrac{n(A)}{n(\Omega)}=\dfrac{1886}{3311}$.}
\end{ex}

% nguồn: Ok_TracNghiemLop11_Lan1
\begin{ex}%[1D9H2-4]%[ma:19fe499f]
	Hộp $A$ có $4$ viên bi trắng, $5$ viên bi đỏ và $6$ viên bi xanh. Hộp $B$ có $7$ viên bi trắng, $6$ viên bi đỏ và $5$ viên bi xanh. Lấy ngẫu nhiên mỗi hộp một viên bi,  tính xác suất để hai viên bi được lấy ra có cùng màu.
	 \choice
		 {$\dfrac{91}{135}$}
		{\True $\dfrac{44}{135}$} 
		{$\dfrac{88}{135}$} 
		{$\dfrac{45}{88}$} 
	 \loigiai{
		Gọi  biến cố
		\begin{itemize}
			\item $A\colon$\lq\lq Hai viên bi  được lấy ra có cùng màu\rq\rq.
			 \item $\mathrm{A}_1\colon$\lq\lq Hai viên bi lấy ra màu trắng\rq\rq. Lúc đó $\mathrm{P}(\mathrm{A}_1)=\dfrac{4}{15}\cdot \dfrac{7}{18}$.
			\item $\mathrm{A}_2\colon$\lq\lq Hai viên bi lấy ra màu đỏ\rq\rq. Lúc  đó $\mathrm{P}(\mathrm{A}_2)=\dfrac{5}{15}\cdot \dfrac{6}{18}$.
			\item $\mathrm{A}_3\colon$\lq\lq Hai viên bi lấy ra  màu xanh\rq\rq. Lúc đó $\mathrm{P}(\mathrm{A}_3)=\dfrac{6}{15}\cdot \dfrac{5}{18}$.
		\end{itemize}
		Lúc đó $A=\mathrm{A}_1\cup{\mathrm{A}_2}\cup{\mathrm{A}_3}$  và $A_1$, $A_2$, $A_3$ là các biến cố xung khắc nên
	$$\mathrm{P}(A)=\mathrm{P}(\mathrm{A}_1)+\mathrm{P}(\mathrm{A}_2)+\mathrm{P}(\mathrm{A}_3)=\dfrac{44}{135}.$$} 
\end{ex}

% nguồn: Ok_TracNghiemLop11_Lan1
\begin{ex}%[1D9H2-4]%[ma:f755fa33]
	Trong  ngân hàng câu hỏi kiểm tra có $20$ câu hỏi trắc nghiệm và $10$ câu hỏi tự luận. Thầy giáo cần chọn ngẫu nhiên $5$ câu hỏi để tạo thành một đề kiểm tra. Xác suất để thầy giáo chọn được đề chỉ có $1$ hoặc $2$ câu hỏi tự luận gần nhất với giá trị nào sau đây?
	\choice 
		 {$0{,}5$}
		{$0{,}6$} 
		 {\True $0{,}7$}
		 {$0{,}8$}
	\loigiai{ 
		 Ta có $n(\Omega)=\mathrm{C}_{30}^5$.\\
		Gọi biến cố $A\colon$\lq\lq Đề có $1$ câu hỏi tự luận\rq\rq.\\ 
		Suy  ra $n(A)=\mathrm{C}_{10}^1\cdot \mathrm{C}_{20}^4$ hay $\mathrm{P}(A)=\dfrac{n(A)}{n(\Omega)}\approx 0{,}34$.\\
		Gọi biến cố $B\colon$\lq\lq Đề có $2$ câu hỏi tự  luận\rq\rq.\\
		 Suy ra $n(B)=\mathrm{C}_{10}^2\cdot \mathrm{C}_{20}^3$ hay $\mathrm{P}(B)=\dfrac{n(B)}{n(\Omega)}\approx 0{,}36$.\\
	Do $A$ và $B$ là hai biến cố  xung khắc nên $\mathrm{P}(A\cup B)=\mathrm{P}(A)+\mathrm{P}(B)\approx 0,7$.}
\end{ex}

% nguồn: Ok_TracNghiemLop11_Lan1
\begin{ex}%[1D9H2-4]%[ma:05b17ab8]
	Có $ 3 $ học sinh lớp $ 12 $A, $ 4 $ học sinh lớp $ 12 $B và $ 5 $ học  sinh lớp $ 12 $C xếp thành hàng dọc. Xác suất để $ 4 $ học sinh lớp $ 12 $B đứng cạnh nhau bằng?
	\choice 
		{\True  $ \dfrac{1}{55} $}
		{$  \dfrac{1}{330} $}
		 {$ \dfrac{2}{11} $}
		 {$ \dfrac{1}{ 27720 } $}
	 \loigiai{
		Ta có  không gian mẫu $ n(\Omega) = 12! $.\\
		Ta  có số biến cố có lợi $ n(A) = 9! \cdot 4! $.\\
		Vậy  xác suất cần tìm $ \mathrm{P}(A) = \dfrac{9! \cdot 4! }{12!} = \dfrac{1}{55} $.
	 }
\end{ex}

% nguồn: Ok_TracNghiemLop11_Lan1
\begin{ex}%[1D9H2-4]%[ma:ab3f4810]
	Cho $A$, $B$ là hai biến cố xung khắc.  Biết $\mathrm{P}(A)=\dfrac{1}{3}$, $\mathrm{P}(B)=\dfrac{1}{4}$. Tính $\mathrm{P}(A\cup B)$.
	 \choice
		{\True  $\dfrac{7}{12}$}
		{$\dfrac{1}{12}$} 
		{$\dfrac{1}{7}$} 
		 {$\dfrac{1}{2}$}
	\loigiai{ 
	$\mathrm{P}(A\cup B)=\mathrm{P}(A)+\mathrm{P}(B)=\dfrac{7}{12}$.} 
\end{ex}

% nguồn: Ok_TracNghiemLop11_Lan1
\begin{ex}%[1D9H2-4]%[ma:e8d6cbba]
	Lớp 11A8 trường THPT X có $25$ học sinh nam và $20$ học sinh nữ. Chọn ngẫu nhiên đồng thời hai bạn từ lớp này để tham dự cuộc họp của trường. Tính xác suất chọn được hai bạn có cùng giới tính để đi dự  cuộc họp.
	\choice 
		{$\dfrac{19}{99}$} 
		{$\dfrac{10}{33}$} 
		{\True $\dfrac{49}{99}$} 
		{$\dfrac{29}{99}$} 
	\loigiai{ 
		Ta  có $n(\Omega)=\mathrm{C}_{45}^2$.\\
		Gọi $A$ là biến cố: ``Chọn được hai bạn nam''.\\ 
		Có  $P(A)=\dfrac{\mathrm{C}_{25}^2}{\mathrm{C}_{45}^2}$.\\
		 Gọi B là biến cố: ``Chọn được hai bạn nữ''.\\
		Có $P(B)=\dfrac{\mathrm{C}_{20}^2}{\mathrm{C}_{45}^2}$.\\ 
		Do $A$ và $B$ là hai biến cố xung khắc nên $P(A  \cup B)=P(A)+P(B)=\dfrac{\mathrm{C}_{25}^2+\mathrm{C}_{20}^2}{\mathrm{C}_{45}^2}=\dfrac{49}{99}$.
	 }
\end{ex}

% nguồn: Ok_TracNghiemLop11_Lan1
\begin{ex}%[1D9H2-4]%[ma:51f4de81]
	Lấy ra ngẫu nhiên $2$ quả bóng từ một hộp chứa $5$ quả bóng xanh và $4$ quả bóng đỏ có kích thước và khối lượng như nhau.  Xác suất của biến cố \lq\lq Hai quả bóng lấy ra có cùng màu\rq\rq\,là
	\choice 
		{$\dfrac{1}{9}$} 
		 {$\dfrac{2}{9}$}
		{\True  $\dfrac{4}{9}$}
		{$\dfrac{5}{9}$} 
	\loigiai{ 
		Số phần tử của  không gian mẫu $n(\Omega)=\mathrm{C}_9^2$.\\
		 Ta gọi các biến cố
		\begin{itemize}
			\item $A$ là biến cố  \lq\lq Hai quả bóng lấy ra đều có màu xanh\rq\rq\,suy ra $\mathrm{P}(A)=\dfrac{\mathrm{C}_5^2}{\mathrm{C}_9^2}$.
			 \item $B$ là biến cố \lq\lq Hai quả bóng lấy ra đều có màu đỏ\rq\rq\,suy ra $\mathrm{P}(B)=\dfrac{\mathrm{C}_4^2}{\mathrm{C}_9^2}$.
			\item $A\cup B$  là biến cố \lq\lq Hai bóng lấy ra có cùng màu\rq\rq.
		\end{itemize}
		Ta thấy $A$ và $B$ xung khắc nên 
	 $$\mathrm{P}(A\cup B)=\mathrm{P}(A)+P(B)=\dfrac{4}{9}.$$}
\end{ex}

% nguồn: Ok_TracNghiemLop11_Lan1
\begin{ex}%[1D9H2-4]%[ma:b0228a6f]
	Trên kệ sách có $8$ sách Toán và $6$ sách Văn khác nhau được xếp thứ tự. Lấy lần lượt $3$ cuốn sách mà không để lại trên kệ. Xác suất để  được ít nhất hai cuốn sách Toán
	\choice 
		 {$\dfrac{8}{14}$}
		{\True $\dfrac{8}{13}$} 
		 {$\dfrac{1}{2}$}
		{$\dfrac{135}{364}$} 
	 \loigiai{
		Số phần tử của không  gian mẫu $n(\Omega)=\mathrm{A}_{14}^3$
		\begin{itemize}
			\item \textit{Trường hợp 1:} Lấy được $3$ cuốn sách Toán  $\mathrm{P}_1=\dfrac{\mathrm{A}_8^3}{\mathrm{A}_{14}^3}=\dfrac{2}{13}$.
			\item \textit{Trường hợp 2:} Lấy được $2$ cuốn sách Toán, $1$  cuốn sách Văn $\mathrm{P}_2=\dfrac{\mathrm{C}_3^2\cdot \mathrm{A}_8^2\cdot \mathrm{A}_6^1}{\mathrm{A}_{14}^3}=\dfrac{6}{13}$.
		\end{itemize}
	Vậy xác suất để được ít nhất  hai cuốn sách Toán là $\mathrm{P}=\mathrm{P}_1+\mathrm{P}_2=\dfrac{8}{13}$.}
\end{ex}

% nguồn: Ok_TracNghiemLop11_Lan1
\begin{ex}%[1D9H2-4]%[ma:36867a00]
	Cho một đồng xu cân đối đồng chất có mặt sấp và mặt ngửa. Gieo đồng xu $6$ lần. Xác suất để số lần xuất hiện mặt ngửa nhiều  hơn số lần xuất hiện mặt sấp là
	\choice 
		 {$\dfrac{21}{32}$}
		{$\dfrac{21}{64}$} 
		{\True $\dfrac{11}{32}$} 
		 {$\dfrac{3}{64}$}
	\loigiai{ 
		Gọi $A$ là biến cố \lq\lq Số lần xuất  hiện mặt ngửa nhiều hơn số lần xuất hiện mặt sấp \rq\rq.\\
		Số phần  tử của không gian mẫu là $n\left(\Omega\right) = 2^6 = 64$.\\
		Ta  có các trường hợp để biến cố $A$ xảy ra sau
		\begin{itemize}
			\item \textbf{Trường hợp 1}: Cả $6$ lần gieo đều là mặt ngửa: Có $1$ cách. 
			\item \textbf{Trường hợp 2}: $5$ ngửa, $1$ sấp: Có  $6$ cách.
			\item \textbf{Trường hợp 3}: $4$ ngửa, $2$ sấp: Có $\dfrac{6!}{4!\cdot 2!}  = 15$ cách.
		\end{itemize}
		Vậy xác  suất cần tìm là $\mathrm P(A) = \dfrac{n(A)}{n\left(\Omega\right)}= \dfrac{1+6+15}{64} = \dfrac{11}{32}$.
	} 
\end{ex}

% nguồn: Ok_TracNghiemLop11_Lan1
\begin{ex}%[1D9H2-4]%[ma:bbe32bee]
	Lớp $11A1$ có $50$ học sinh, trong đó có $32$ bạn thích học môn Toán, $17$ bạn thích học môn Hóa và $8$ bạn thích cả hai môn trên. Chọn  ngẫu nhiên một bạn. Xác suất để bạn đó thích môn Toán hoặc Hóa là?
	 \choice
		{$\dfrac{49}{50}$} 
		 {$\dfrac{9}{50}$}
		 {$\dfrac{33}{50}$}
		 {\True $\dfrac{41}{50}$}
	\loigiai{ 
		Biến cố  \lq\lq Bạn đó thích môn Toán hoặc Hóa\rq\rq\, là $A\cup B$.\\
		Theo công thức cộng xác suất  ta có $\mathrm{P}\left(A\cup B\right)=\mathrm{P}(A)+\mathrm{P}(B)-\mathrm{P}\left(A\cap B\right) = \dfrac{32}{50}+\dfrac{17}{50}-\dfrac{8}{50}=\dfrac{41}{50}$.
	} 
\end{ex}


%%%%%%%%%%%%%%%% ĐÚNG/SAI %%%%%%%%%%%%%%%%

% nguồn: Ok_dungsai_lop11_Lan1
\begin{ex}%[1D9H2-4]%[ma:8cc750b3]
 Túi $X$ chứa ba viên bi trắng và hai viên bi đỏ. Túi $Y$ chứa một bi trắng và ba màu đỏ viên bi. Người ta chọn ngẫu nhiên mỗi hộp một viên bi. Gọi $A$ là biến cố \lq\lq Lấy được viên bi màu trắng từ  túi $X$\rq\rq. Gọi $B$ là biến cố \lq\lq Lấy được viên bi màu trắng từ túi $Y$\rq\rq. Gọi $X_1$ là biến cố \lq\lq Lấy được hai viên bi cùng màu đỏ\rq\rq.
  \choiceTF
   {\True $\mathrm{P}(A)=\dfrac{3}{5}$}
   {$\mathrm{P}(B)=\dfrac{1}{5}$}
  {$\mathrm{P}(X_1)=\dfrac{4}{5}$} 
  {\True Xác  suất để lấy được hai viên bi cùng màu bằng $\mathrm{P}(X)=\dfrac{7}{15}$}
  \loigiai{
  \begin{itemchoice}
   \itemch Vì hai túi $X$, $Y$ đọc lập nên ta  chỉ cần xét riêng túi $X$.\\
    Số cách để lấy được $1$ viên bi trắng trong túi $X$ là $3$. Tổng số viên bi của túi $X$ là $5$.\\
   Suy ra $\mathrm{P}(A)=\dfrac{3}{5}$. 
   \itemch Vì hai túi $X$,  $Y$ đọc lập nên ta chỉ cần xét riêng túi $Y$.\\
    Số cách để lấy được $1$ viên bi trắng trong túi $Y$ là $1$. Tổng số viên bi của túi $Y$ là $4$.\\
   Suy ra $\mathrm{P}(B)=\dfrac{1}{4}$. 
   \itemch Vì $A$ và  $B$ là hai biến cố độc lập và $X_1=A\cap B$ nên $\mathrm{P}(X_1)=\mathrm{P}(A)\cdot \mathrm{P}(B)=\dfrac{3}{5}\cdot\dfrac{1}{3}=\dfrac{1}{5}$.
   \itemch Gọi $X_1$ là  biến cố \lq\lq Lấy được hai viên bi cùng màu đỏ\rq\rq.\\
   Vì $\overline{A}$ và $\overline{B}$ là hai biến cố độc lập và $X_1=\overline{A}\cap\overline{B}$ nên  $\mathrm{P}(X_1)=\mathrm{P}(\overline{A})\cdot \mathrm{P}(\overline{B})=\dfrac{2}{5}\cdot\dfrac{2}{3}=\dfrac{4}{15}$.\\
   Biến cố để hai viên bi lấy ra cùng màu  là $X=X_1\cup X_2$\\
   Lại có $X_1$ và $X_2$ là hai biến cố xung khắc, xác suất để hai viên bi lấy ra cùng  màu là\\
   $\mathrm{P}(X)=\mathrm{P}(X_1)+\mathrm{P}(X_2)=\dfrac{4}{15}+\dfrac{1}{5}=\dfrac{7}{15}$. 
  \end{itemchoice}
  }
\end{ex}

% nguồn: Ok_dungsai_lop11_Lan1
\begin{ex}%[1D9H2-4]%[ma:5d9ef263]
 Trong một công ty  có $40$ nhân viên, trong đó có $27$ người thích chơi bóng bàn, $25$ người thích chơi cầu lông. Chọn ngẫu nhiên một nhân viên trong công ty đó. Các khẳng định sau đây khẳng định nào đúng, khẳng định nào sai?
 \choiceTF 
  {\True Xác suất để người đó thích chơi ít nhất một trong hai môn bóng bàn và cầu lông  không vượt quá $100\%$}
  {Xác suất để người  đó thích chơi cả hai môn cầu lông và bóng bàn ít nhất là $25\%$}
  {Xác suất để người đó không thích chơi bóng bàn và không thích chơi cầu  lông ít nhất là $2{,}5\%$}
  {Xác suất  để người đó không thích chơi bóng bàn và không thích chơi cầu lông nhiều nhất là $25\%$}
  \loigiai{
  Gọi $A$ là biến cố  \lq\lq Người đó thích chơi bóng bàn\rq\rq;\\
  $B$ là biến cố \lq\lq  Người đó thích chơi cầu lông\rq\rq.\\
  Theo bài  ra ta có $n(A)=27,n(B)=25$.\\
  Gọi $x=n(A\overline{B}),y=n(AB),z=n(\overline{A}B),t=n(\overline{A}\,\overline{B})$. Khi đó  ta có
  $$\heva{& x+y=27\\& y+z=25\\ & x+y+z+t=40} \Rightarrow \heva{& y-t=12\\ & x+y=27\\ & y+z=25} \Rightarrow \heva{& 12\le y\le 25\\ & 0\le t\le  25-12=13} \Rightarrow \heva{& \dfrac{12}{40}\le \mathrm{P}(AB)\le \dfrac{25}{40} \\ & 0\le \mathrm{P}(\overline{A}\,\overline{B})\le \dfrac{13}{40}.}$$
  \begin{itemchoice}
   \itemch Biến cố đối của biến cố $A\cup B\colon$ \lq\lq Người đó thích chơi ít nhất một trong hai môn bóng bàn và cầu lông\rq\rq~ là biến cố $\overline{A}\,\overline{B}\colon$ \lq\lq Người đó không thích chơi  cầu lông và bóng bàn\rq\rq.\\
   Ta có  $\mathrm{P}(A)=\dfrac{27}{40};\mathrm{P}(B)=\dfrac{25}{40} \Rightarrow \mathrm{P}(A\cup B)=1-\mathrm{P}(\overline{A}\,\overline{B})\le 1$.
   \itemch Ta  cần tính $\mathrm{P}(AB)$.\\
   Ta có $\mathrm{P}(AB)=\mathrm{P}(A)+\mathrm{P}(B)-\mathrm{P}(A\cup  B)=\dfrac{27}{40}+\dfrac{25}{40}-\left[1-\mathrm{P}(\overline{A}\,\overline{B}) \right]=\dfrac{12}{40}+\mathrm{P}(\overline{A}\,\overline{B})\ge \dfrac{12}{40}=30\%$.
   \itemch  Ta cần tính $\mathrm{P}(\overline{A}\,\overline{B})$.\\
   Ta có $\mathrm{P}(\overline{A}\,\overline{B})=1-\mathrm{P}(A\cup B)=1-\mathrm{P}(A)-\mathrm{P}(B)+\mathrm{P}(AB)=\mathrm{P}(AB)-\dfrac{12}{40}\ge 0$. 
   \itemch Ta  cần tính $\mathrm{P}(\overline{A}\,\overline{B})$.\\
   Ta  có $\mathrm{P}(\overline{A}\,\overline{B})=1-\mathrm{P}(A\cup B)=1-\mathrm{P}(A)-\mathrm{P}(B)+\mathrm{P}(AB)=\mathrm{P}(AB)-\dfrac{12}{40}\le \dfrac{25}{40}-\dfrac{12}{40}=\dfrac{13}{40}=32{,}5\%$.
  \end{itemchoice}
 } 
\end{ex}

% nguồn: Ok_dungsai_lop11_Lan1
\begin{ex}%[1D9H2-4]%[ma:00be8452]
 Chọn ngẫu nhiên một số nguyên dương có hai chữ số. Xét biến cố $A$: \lq\lq Số được viết ra là số chia hết cho $8$\rq\rq \,  và biến cố $B$: \lq\lq Số được viết ra là số chia hết cho $9$\rq\rq.
 \choiceTF 
  {Số  phần tử của không gian mẫu $n(\Omega)=89$}
  {\True Xác xuất của biến cố  $A$ là $\mathrm{P}(A)=\dfrac{11}{90}$}
  {\True Xác xuất của biến cố $B$  là $\mathrm{P}(B)=\dfrac{1}{9}$}
  {Xác xuất của biến cố hợp là $\mathrm{P}(A  \cup B)=\dfrac{1}{10}$}
 \loigiai{ 
  Trong $90$ số có  hai chữ số thì
  \begin{itemize}
   \item có $11$ số  chia hết cho $8$,
   \item  có $10$ số chia hết cho $9$,
   \item có $1$ số chia hết cho cả $8$ và  $9$.
  \end{itemize}
  \begin{itemchoice}
    \itemch Ta có $n(\Omega)=90$.
   \itemch Số các  kết quả thuận lợi cho các biến cố $A$ là $n(A)=11$.\\
   Suy ra $\mathrm{P}(A)=\dfrac{n(A)}{n(\Omega)}=\dfrac{11}{90}$. 
   \itemch Số các  kết quả thuận lợi cho các biến cố $B$ à $n(B)=9$.\\
   Suy ra $\mathrm{P}(B)=\dfrac{n(B)}{n(\Omega)}=\dfrac{9}{90}=\dfrac{1}{10}$. 
    \itemch Ta có $\mathrm{P}(A)=\dfrac{11}{90}$, $\mathrm{P}(B)=\dfrac{10}{90}$, $\mathrm{P}(A \cap B)=\dfrac{1}{90}$.\\
    Vậy $\mathrm{P}(A \cup B)=\mathrm{P}(A)+\mathrm{P}(B) - \mathrm{P}(A \cap B) =\dfrac{11}{90}+\dfrac{10}{90}-\dfrac{1}{90}=\dfrac{2}{9}$.
  \end{itemchoice}
  }
\end{ex}

% nguồn: Ok_dungsai_lop11_Lan1
\begin{ex}%[1D9H2-4]%[ma:ab48e3d9]
 Một lớp học có $38$ học sinh, trong đó có $18$ học sinh thích chơi bóng đá, $16$ học sinh thích chơi cầu lông, $10$  học sinh không thích chơi cả bóng đá và cầu lông. Chọn ngẫu nhiên một học sinh trong lớp đó. Trong các khẳng định sau, khẳng định nào đúng, khẳng định nào sai?
 \choiceTF 
  {\True Xác suất để học sinh đó  thích chơi cầu lông bằng $\dfrac{8}{19}$}
  {Xác suất để học sinh đó thích chơi ít nhất một trong hai môn bóng đá và cầu  lông bằng $\dfrac{5}{19}$}
  {Xác suất để học sinh đó thích chơi cầu lông và không thích chơi  bóng đá bằng $\dfrac{6}{19}$}
  {\True Thích  chơi đúng một trong hai môn bằng $\dfrac{11}{19}$}
  \loigiai{
  Gọi $A$ là  biến cố \lq\lq Học sinh đó thích chơi bóng đá\rq\rq;\\
  $B$ là  biến cố \lq\lq Học sinh đó thích chơi cầu lông\rq\rq.
  \begin{itemchoice}
   \itemch Ta có $n(B)=16,\, n(\Omega)=38 \Rightarrow  \mathrm{P}(B)=\dfrac{8}{19}$.
   \itemch  Ta cần tính $\mathrm{P}(A\cup B)$.\\
   Biến cố đối của biến cố $A\cup B\colon$ \lq\lq Học sinh đó thích chơi ít nhất một  trong hai môn\rq\rq~ là biến cố $\overline{AB}\colon$ \lq\lq Học sinh đó không thích chơi cả hai môn cầu lông và bóng đá\rq\rq.\\
   Ta có $\mathrm{P}(\overline{AB})=\dfrac{5}{19} \Rightarrow \mathrm{P}(A\cup  B)=1-\mathrm{P}(\overline{AB})=1-\dfrac{5}{19}=\dfrac{14}{19}$.
   \itemch Ta cần tính  $\mathrm{P}(\overline{A}B)$.\\
   Ta  có $\mathrm{P}(A)=\dfrac{9}{19}$, $\mathrm{P}(B)=\dfrac{8}{19} \Rightarrow \mathrm{P}(AB)=\mathrm{P}(A)+\mathrm{P}(B)-\mathrm{P}(A\cup B)=\dfrac{9}{19}+\dfrac{8}{19}-\dfrac{14}{19}=\dfrac{3}{19}$.\\
   Lại có $B=AB \cup \overline{A}B \Rightarrow \mathrm{P}(B)=\mathrm{P}(AB)+\mathrm{P}(\overline{A}B)  \Rightarrow \mathrm{P}(\overline{A}B)=\mathrm{P}(B)-\mathrm{P}(AB)=\dfrac{8}{19}-\dfrac{3}{19}=\dfrac{5}{19}$.
   \itemch  Gọi $E$ là biến cố \lq\lq Học sinh đó thích chơi đúng một trong hai môn cầu lông hay bóng đá\rq\rq.\\
   Ta có $A=AB\cup A\overline{B} \Rightarrow \mathrm{P}(A)=\mathrm{P}(AB)+\mathrm{P}(A\overline{B}) \Rightarrow \mathrm{P}(A\overline{B})=\mathrm{P}(A)-\mathrm{P}(AB)=\dfrac{9}{19}-\dfrac{3}{19}=\dfrac{6}{19}$.\\ 
   Ta  có $E=A\overline{B}\cup \overline{A}B$, suy ra $\mathrm{P}(E)=\mathrm{P}(A\overline{B})+\mathrm{P}(\overline{A}B)=\dfrac{6}{19}+\dfrac{5}{19}=\dfrac{11}{19}$.
  \end{itemchoice}
  }
\end{ex}

% nguồn: Ok_dungsai_lop11_Lan1
\begin{ex}%[1D9H2-4]%[ma:a09c28f9]
 Một hộp đựng $4$ viên bi màu xanh, $3$ viên bi màu đỏ và $2$ viên bi màu vàng. Chọn ngẫu  nhiên $2$ viên bi từ hộp trên. Gọi $A$ là biến cố \lq\lq Chọn được $2$ viên bi màu xanh\rq\rq, $B$ là biến cố \lq\lq Chọn được $2$ viên bi màu đỏ\rq\rq, $C$ là biến cố \lq\lq Chọn được $2$ viên bi màu vàng\rq\rq. Khi đó
 \choiceTF 
   {$\mathrm{P}(A)=\dfrac{1}{7}$}
  {$\mathrm{P}(B)=\dfrac{1}{8}$} 
  {\True $\mathrm{P}(C)=\dfrac{1}{36}$} 
  {\True Xác suất để  chọn được $2$ viên bi cùng màu bằng $\mathrm{P}(X)=\mathrm{P}(A)+\mathrm{P}(B)+\mathrm{P}(C)=\dfrac{1}{6}+\dfrac{1}{12}+\dfrac{1}{36}=\dfrac{5}{18}$}
 \loigiai{ 
  Số phần tử của không gian  mẫu $n(\Omega)=\mathrm{C}_9^2$.
  Ta  có
  \begin{itemchoice}
   \itemch $\mathrm{P}(A)=\dfrac{\mathrm{C}_4^2}{\mathrm{C}_9^2}=\dfrac{1}{6}$; 
   \itemch  $ \mathrm{P}(B)=\dfrac{\mathrm{C}_3^2}{\mathrm{C}_9^2}=\dfrac{1}{12}$;
   \itemch $  \mathrm{P}(C)=\dfrac{\mathrm{C}_2^2}{\mathrm{C}_9^2}=\dfrac{1}{36}$.
   \itemch Ta có $X=A\cup B\cup C$ và các biến cố $A$, $B$, $C$ đôi một xung khắc.\\ 
   Do đó, ta có  $\mathrm{P}(X)=\mathrm{P}(A)+\mathrm{P}(B)+\mathrm{P}(C)=\dfrac{1}{6}+\dfrac{1}{12}+\dfrac{1}{36}=\dfrac{5}{18}$.
  \end{itemchoice}
 } 
\end{ex}

% nguồn: Ok_dungsai_lop11_Lan1
\begin{ex}%[1D9H2-4]%[ma:8f8d7669]
 Lớp $11$A của một trường có $30$ học sinh, trong đó có $12$ bạn thích nhạc cổ điển, $11$ bạn thích nhạc trẻ và $3$ bạn thích cả nhạc cổ điển và nhạc trẻ. Chọn ngẫu nhiên một bạn trong  lớp. Các khẳng định sau đây, khẳng định nào là đúng, khẳng định nào là sai?
 \choiceTF 
  {Biến cố bạn đó thích nhạc cổ điển và biến cố bạn đó thích nhạc trẻ là hai biến  cố xung khắc}
   {Xác suất để bạn đó thích nhạc trẻ là $\dfrac{2}{5}$, xác suất để bạn đó thích nhạc cổ điển là $\dfrac{11}{30}$}
  {\True Xác suất để bạn đó thích nhạc cổ điển hoặc nhạc  trẻ là $\dfrac{2}{3}$}
  {\True Xác suất để bạn đó không thích nhạc cổ điển và không thích  nhạc trẻ là $\dfrac{1}{3}$}
  \loigiai{
  Gọi $A$ là biến  cố \lq\lq Bạn đó thích nhạc cổ điển\rq\rq;\\
  $B$ là biến cố \lq\lq Bạn đó thích nhạc  trẻ\rq\rq;\
  $AB$ là biến cố \lq\lq Bạn đó thích cả nhạc cổ  điển và nhạc trẻ\rq\rq.
  \begin{itemchoice}
   \itemch Vì có $3$ bạn thích cả nhạc cổ  điển và nhạc trẻ nên biến cố $A$ và biến cố $B$ không là hai biến cố xung khắc.
   \itemch Theo bài ra ta có $\mathrm{P}\left(A \right)=\dfrac{12}{30}=\dfrac{2}{5};\mathrm{P}\left(B \right)=\dfrac{11}{30};\mathrm{P}\left(AB \right)=\dfrac{3}{30}=\dfrac{1}{10}$. 
    \itemch Ta có $A\cup B$ là biến cố \lq\lq Bạn đó thích nhạc cổ điển hoặc nhạc trẻ\rq\rq.\\
   Áp dụng công thức cộng  xác suất ta có $\mathrm{P}\left(A\cup B \right)=\mathrm{P}\left(A \right)+\mathrm{P}\left(B \right)-\mathrm{P}\left(AB \right)=\dfrac{12}{30}+\dfrac{11}{30}-\dfrac{3}{30}=\dfrac{2}{3}$.
   \itemch Ta có $\overline{AB}$ là biến cố \lq\lq Bạn đó không thích cả nhạc cổ điển và nhạc  trẻ\rq\rq.\\
   Vì $A\cup B$  là biến cố đối của biến cố $\overline{AB}$ nên $\mathrm{P}\left(\overline{AB} \right)=1-\mathrm{P}\left(A\cup B \right)=\dfrac{1}{3}$.
  \end{itemchoice}
  }
\end{ex}

% nguồn: Ok_dungsai_lop11_Lan1
\begin{ex}%[1D9H2-4]%[ma:c7f44167]
 Một hộp đựng $10$ tấm thẻ được đánh số từ $1$ đến $10$, hai tấm thẻ khác nhau đánh hai số  khác nhau. Rút ngẫu nhiên một tấm thẻ, khi đó
 \choiceTF 
   {\True Gọi $A$ là biến cố \lq\lq Rút được thẻ đánh số chia hết cho $2$\rq\rq, suy ra $n(A)=5$}
  {\True Gọi $A$ là biến cố \lq\lq Rút được thẻ đánh số chia hết cho $2$\rq\rq,  suy ra $\mathrm{P}(A)=\dfrac{1}{2}$}
  {Gọi $B$ là biến cố  \lq\lq Rút được thẻ đánh số chia hết cho $7$\rq\rq, suy ra $\mathrm{P}(B)=\dfrac{1}{8}$}
   {Xác suất để rút được thẻ đánh số chia hết cho $2$ hoặc $7$ bằng $\dfrac{3}{7}$}
 \loigiai{ 
  \begin{itemchoice}
    \itemch $A=\{2; 4; 6; 8; 10\}$ suy ra $n(A)=5$.
    \itemch $\mathrm{P}(A)=\dfrac{5}{10}=\dfrac{1}{2}$.
   \itemch $B=\{7\}$ và  $\mathrm{P}(B)=\dfrac{1}{10}$.
   \itemch Vì $A$  và $B$ là hai biến cố xung khắc nên $\mathrm{P}(A\cup B)=\mathrm{P}(A)+\mathrm{P}(B)=\dfrac{1}{2}+\dfrac{1}{10}=\dfrac{3}{5}$.
  \end{itemchoice}
  }
\end{ex}

% nguồn: Ok_dungsai_lop11_Lan1
\begin{ex}%[1D9H2-4]%[ma:a72589c2]
 Mai, Lan và $5$ bạn cùng lớp xếp thành một hàng ngang theo thứ tự ngẫu nhiên. Gọi $A$ là  biến cố \lq\lq Mai đứng ở đầu hàng\rq\rq\,, $B$ là biến cố \lq\lq Lan đứng ở đầu hàng\rq\rq. Khi đó
  \choiceTF
   {$\mathrm{P}(A)=\dfrac{1}{7}$}
   {\True $\mathrm{P}(B)=\dfrac{2}{7}$}
  {\True  $\mathrm{P}(AB)=\dfrac{21}{21}$}
  {$\mathrm{P}(A \cup B)=\dfrac{10}{21}$} 
  \loigiai{
  Số  cách xếp $7$ người thành một hàng ngang là $7!$.
  \begin{itemchoice}
   \itemch Xác suất của biến  cố $A$ là $\mathrm{P}(A)=\dfrac{2 \cdot 6 !}{7 !}=\dfrac{2}{7}$.
   \itemch Xác suất của  biến cố $B$ là $\mathrm{P}(B)=\dfrac{2 \cdot 6 !}{7 !}=\dfrac{2}{7}$.
   \itemch Xác suất của biến cố \lq\lq  Hai bạn Lan và Mai đứng ở hai đầu hàng\rq\rq~là
   $$ 
   \mathrm{P}(A B)=\dfrac{2 \cdot 5 !}{7  !}=\dfrac{1}{21}.
   $$ 
   \itemch Xác suất của biến cố \lq\lq Có ít nhất một trong hai bạn Mai và Lan đứng ở đầu  hàng\rq\rq~là
    $$
   \mathrm{P}(A  \cup B)=\mathrm{P}(A)+\mathrm{P}(B)-\mathrm{P}(A B)=\dfrac{2}{7}+\dfrac{2}{7}-\dfrac{1}{21}=\dfrac{11}{21}.
    $$
  \end{itemchoice}
  }
\end{ex}

% nguồn: Ok_dungsai_lop11_Lan1
\begin{ex}%[1D9H2-4]%[ma:de490e0e]
 Lớp 11A1 có $50$ học sinh, trong đó có $32$ bạn thích học môn Toán, $17$ bạn thích học môn Lịch Sử và $8$ bạn thích  cả hai môn trên. Chọn ngẫu nhiên một bạn trong lớp. Các khẳng định sau đây, khẳng định nào đúng, khẳng định nào sai?
 \choiceTF 
   {\True Xác suất để bạn đó thích học môn Toán là $\dfrac{16}{25}$}
  {\True Xác suất để bạn đó thích cả môn Toán và môn Lịch  Sử là $\dfrac{4}{25}$}
  {\True Xác suất  để bạn đó thích môn Toán hoặc môn Lịch Sử là $\dfrac{41}{50}$}
  {Xác suất để bạn đó  không thích cả môn Toán và môn Lịch Sử là $\dfrac{1}{50}$}
  \loigiai{
  Gọi $A$  là biến cố \lq\lq Học sinh thích học môn Toán\rq\rq;\\
  $B$ là biến  cố \lq\lq Học sinh thích học môn Lịch sử\rq\rq.
  \begin{itemchoice}
   \itemch Ta có $n(A)=32$ và  $n(\Omega)=50$ nên $\mathrm{P}(A)=\dfrac{32}{50}=\dfrac{16}{25}$.
   \itemch $A\cap B$ là biến cố \lq\lq Học sinh thích cả môn Toán và môn Lịch Sử\rq\rq.\\ 
   Ta  có $n(A\cap B)=8$ và $n(\Omega)=50$ nên $\mathrm{P}(A\cap B)=\dfrac{8}{50}=\dfrac{4}{25}$.
   \itemch $A\cup B$ là biến cố \lq\lq Học sinh thích  môn Toán hoặc môn Lịch Sử\rq\rq.\\
   Theo công thức cộng xác suất, ta  có $\mathrm{P}(A\cup B)=\mathrm{P}(A)+\mathrm{P}(B)-\mathrm{P}(A\cap B)=\dfrac{16}{25}+\dfrac{17}{50}-\dfrac{4}{25}=\dfrac{41}{50}$.
   \itemch $\overline{A\cup B}$ là biến cố \lq\lq Học sinh không thích cả môn Toán  và môn Lịch Sử\rq\rq.\\
   Theo tính chất  xác suất, ta có $\mathrm{P}(\overline{A\cup B})=1-\mathrm{P}(A\cup B)=1-\dfrac{41}{50}=\dfrac{9}{50}$.
  \end{itemchoice}
  }
\end{ex}

% nguồn: Ok_dungsai_lop11_Lan1
\begin{ex}%[1D9C2-4]%[ma:56a136dc]
 Một nhà sách bán hai loại  sách $I$ và $II$. Theo thống kê cho thấy có $60\%$ số người mua loại sách $I$, $70\%$ số người mua loại sách $II$ và có $50\%$ số người mua cả hai loại sách trên. Chọn ngẫu nhiên một người mua. Xét biến cố $A\colon$\lq\lq Người đó mua loại sách $I$\rq\rq và $B$ là biến cố$\colon$\lq\lq Người đó mua loại sách $II$\rq\rq. Xét tính Đúng - Sai của các khẳng định sau:
  \choiceTF
  {\True $\mathrm{P}(A)=0{,}6$  và $\mathrm{P}(B)=0{,}7$}
  {\lq\lq Người đó  mua ít nhất một trong hai loại sách\rq\rq\, là biến cố $A\cap B$}
  {Xác suất người đó mua ít nhất một trong hai loại sách  là $0{,}9$}
  {\True Xác suất người đó  không mua hai loại sách trên là $0{,}2$}
 \loigiai{ 
  \begin{itemchoice}
   \itemch  Đúng.
   \itemch Sai.\\ 
   \lq\lq Người đó mua ít nhất  một trong hai loại sách\rq\rq\, là biến cố $A\cup B$.
    \itemch Sai.\\
   Ta có biến cố \lq\lq Người  đó mua ít nhất một trong hai loại sách\rq\rq\, là $A\cup B$.\\
   Theo đề bài ta có $\mathrm{P}\left(A\cap B\right)=0{,}5$.\\ 
   $\mathrm{P}\left(A\cup B\right)=\mathrm{P}(A)+\mathrm{P}(B)-\mathrm{P}\left(A\cap  B\right) = 0{,}6+0{,}7-0{,}5=0{,}8$.\\
   Vậy xác suất người đó mua  ít nhất một trong hai loại sách là $0{,}8$
    \itemch Đúng.\\
   Gọi $C$ là biến cố$\colon$ \lq\lq Người đó không mua  cả hai loại sách $I$ và $II$\rq\rq.\\
   Khi đó biến cố $\overline{C}=A\cup B \Rightarrow \mathrm{P}\left(\overline{C}\right)=0{,}8 \Rightarrow \mathrm{P}(C)=0{,}2$.\\ 
    Vậy xác suất người đó không mua hai loại sách trên là $0{,}2$.
  \end{itemchoice}
  }
\end{ex}


%%%%%%%%%%%%%%%% TRẢ LỜI NGẮN %%%%%%%%%%%%%%%%

% nguồn: Ok_TraLoiNgan_Lop11_Lan1
\begin{ex}%[1D9N2-4]%[ma:cb904579]
	Một hộp có chứa một số quả cầu gồm bốn màu xanh, vàng, đỏ, trắng (các quả cầu cùng màu thì khác nhau về bán kính). Lấy ngẫu nhiên một quả cầu từ hộp, biết xác suất để lấy được một quả cầu màu xanh bằng $\dfrac{1}{4}$, xác suất để lấy được một quả cầu màu vàng bằng $\dfrac{1}{3}$. Tính xác suất để lấy được một quả cầu xanh hoặc một quả cầu vàng(làm  tròn đến hàng phần trăm).
	\shortans{0,58} 
	\loigiai{ 
		 Gọi biến cố $A\colon$ \lq\lq Lấy được một quả cầu màu xanh\rq\rq\, và $B\colon$ \lq\lq Lấy được một quả cầu màu vàng\rq\rq.\\
		Ta có $A$, $B$ là hai biến cố xung khắc. Suy ra xác suất để lấy được một quả cầu màu  xanh hoặc một quả cầu màu vàng là
	$$\mathrm{P}(A\cup  B)=\mathrm{P}(A)+\mathrm{P}(B)=\dfrac{1}{4}+\dfrac{1}{3}=\dfrac{7}{12}.$$}
\end{ex}

% nguồn: Ok_TraLoiNgan_Lop11_Lan1
\begin{ex}%[1D9N2-4]%[ma:557648a8]
	Chọn ngẫu nhiên $2$  đỉnh trong số $20$ đỉnh của một đa giác đều $20$ cạnh. Tính xác suất của biến cố $A$ \lq\lq$2$ đỉnh được chọn là đường chéo của đa giác\rq\rq\,(làm tròn đến hàng phần trăm).
	\shortans{0,89} 
	 \loigiai{
		Đường nối $2$ đỉnh bất kỳ của một đa giác đều hoặc là cạnh của  đa giác hoặc là đường chéo của đa giác. Số đường chéo của một đa giác đều $20$ cạnh là $\mathrm{C}_{20}^2-20=170$.\\
	Xác suất  của biến cố $A$ là $\mathrm{P}(A)=\dfrac{\mathrm{C}_{20}^2-20}{\mathrm{C}_{20}^2}=\dfrac{17}{19}$.}
\end{ex}

% nguồn: Ok_TraLoiNgan_Lop11_Lan1
\begin{ex}%[1D9H2-4]%[ma:67a08b3b]
	Một lớp học gồm $50$ bạn, trong đó có $20$ bạn thích chơi bóng đá, $28$ bạn thích chơi bóng  rổ và $8$ bạn thích chơi cả hai môn. Gặp ngẫu nhiên $1$ học sinh trong lớp. Xác suất của biến cố \lq\lq Bạn được gặp thích chơi bóng đá hoặc bóng rổ\rq\rq\, là
	\shortans{0,8} 
	\loigiai{ 
		Gọi  $A$ là biến cố \lq\lq Bạn được gặp thích chơi bóng đá\rq\rq,
		 $B$ là biến cố \lq\lq Bạn được gặp thích chơi bóng rổ\rq\rq.\\
		 Xác suất của biến cố $A$ là $$\mathrm{P}(A)=\dfrac{20}{50}=\dfrac{2}{5}.$$
		Xác  suất của biến cố $B$ là $$\mathrm{P}(B)=\dfrac{28}{50}=\dfrac{14}{25}.$$
		Xác suất  của biến cố \lq\lq Bạn được gặp thích chơi cả hai loại\rq\rq \ là:
		$$\mathrm{P}(A B)=\dfrac{8}{50}=\dfrac{4}{25}.$$ 
		Xác suất của biến cố  \lq\lq Bạn được gặp thích chơi bóng đá hoặc bóng rổ \rq\rq\, là:
		$$\mathrm{P}(A  \cup B)=\mathrm{P}(A)+\mathrm{P}(B)-\mathrm{P}(AB)=\dfrac{2}{5}+\dfrac{14}{25}-\dfrac{4}{25}=\dfrac{4}{5}=0{,}8.$$
	 }
\end{ex}

% nguồn: Ok_TraLoiNgan_Lop11_Lan1
\begin{ex}%[1D9H2-4]%[ma:ee156904]
	Một đội tình nguyện gồm $6$ học sinh khối $11$, và $8$ học sinh khối $12$. Chọn ra ngẫu nhiên $2$ người trong đội. Tính xác suất của biến cố  \lq\lq Cả hai người được chọn học cùng một khối\rq\rq\,(làm tròn đến hàng phần trăm).
	\shortans{0,47} 
	 \loigiai{
		Gọi $A$ là biến cố \lq\lq Cả hai học sinh được chọn đều thuộc khối $11$\rq\rq. Gọi $B$ là biến cố \lq\lq Cả hai học sinh được chọn đều thuộc khối $12$\rq\rq. Khi đó $A\cup B$ là biến cố \lq\lq Cả hai  người được chọn học cùng một khối\rq\rq.\\
	Do đó $A$ và $B$  là hai biến cố xung khắc nên $\mathrm{P}(A\cup B)=\mathrm{P}(A)+\mathrm{P}(B)=\dfrac{\mathrm{C}_6^2}{\mathrm{C}_{14}^2}+\dfrac{\mathrm{C}_8^2}{\mathrm{C}_{14}^2}=\dfrac{43}{91}$.}
\end{ex}

% nguồn: Ok_TraLoiNgan_Lop11_Lan1
\begin{ex}%[1D9H2-4]%[ma:49c52f8f]
	Một hộp đựng nhiều quả cầu với nhiều màu sắc  khác nhau. Người ta lấy ngẫu nhiên một quả cầu từ hộp đó. Biết xác suất để lấy được một quả cầu màu xanh từ hộp bằng $\dfrac{1}{5}$, xác suất để lấy được một quả cầu màu đỏ từ hộp bằng $\dfrac{1}{6}$. Gọi $A$ là biến cố \lq\lq Lấy được một quả cầu màu xanh\rq\rq\,và $B$ là biến cố \lq\lq Lấy được một quả cầu màu đỏ\rq\rq. Tính xác suất để lấy được một quả cầu màu xanh hoặc một quả cầu màu đỏ từ hộp (làm tròn đến hàng phần trăm).
	\shortans{0,37} 
	\loigiai{ 
		Mỗi lần lấy thì ta chỉ lấy có một quả cầu, nên nếu lấy được quả cầu  màu xanh thì không có quả cầu màu đỏ (và ngược lại), nói cách khác $\mathrm{P}(A B)=0$.\\
		Vì vậy $A$, $B$ là hai  biến cố xung khắc.\\
		Xác suất để lấy được một quả cầu màu xanh  hoặc một quả cầu màu đỏ là
	$$\mathrm{P}(A\cup  B)=\mathrm{P}(A)+\mathrm{P}(B)=\dfrac{1}{5}+\dfrac{1}{6}=\dfrac{11}{30}.$$}
\end{ex}

% nguồn: Ok_TraLoiNgan_Lop11_Lan1
\begin{ex}%[1D9H2-4]%[ma:2e3a446d]
	Hai bạn Chiến và Công cùng chơi cờ với nhau. Trong một ván cờ, xác suất Chiến thắng Công là $0{,}3$ và xác suất để Công thắng Chiến là $0{,}4$. Hai bạn dừng chơi khi có người thắng, người  thua. Tính xác suất để hai bạn dừng chơi sau hai ván cờ.
	\shortans{0,21} 
	 \loigiai{
		Gọi 
		\begin{itemize}
			\item $A$ là biến  cố \lq\lq Chiến thắng Công trong ván cờ\rq\rq;
			\item $B$  là biến cố \lq\lq Công thắng Chiến trong ván cờ\rq\rq;
			\item $C$ là biến cố \lq\lq Công và Chiến hoà nhau trong ván cờ\rq\rq. 
		\end{itemize}
		Dễ thấy $A$, $B$, $C$ là các biến  cố xung khắc.\\
		Theo giả thiết thì ván đấu thứ nhất hai bạn hoà nhau,  ván đấu thứ hai sẽ có thắng thua.
		\begin{itemize}
			 \item Xét ván thứ nhất: $\mathrm{P}(C)=1-\mathrm{P}(A)-\mathrm{P}(B)=1-0{,}3-0{,}4=0{,}3$.
			\item Xét  ván thứ hai: $\mathrm{P}(A\cup B)=\mathrm{P}(A)+\mathrm{P}(B)=0{,}3+0{,}4=0{,}7$.
		\end{itemize}
	Xác suất để hai bạn  dừng chơi sau hai ván đấu là $P=0{,}3\cdot 0{,}7=0{,}21$.}
\end{ex}

% nguồn: Ok_TraLoiNgan_Lop11_Lan1
\begin{ex}%[1D9H2-4]%[ma:8243a99f]
	Một hộp đựng $9$ tấm thẻ được đánh số từ $1$ tới $9$, hai tấm thẻ khác nhau đánh hai  số khác nhau. Rút ngẫu nhiên đồng thời hai tấm thẻ từ hộp. Xét các biến cố sau\\
	$A\colon$ \lq\lq Cả hai tấm thẻ đều đánh số chẵn\rq\rq, $B\colon$ \lq\lq Chỉ có một tấm thẻ đánh số chẵn\rq\rq, $C\colon$ \lq\lq Tích hai số đánh trên hai tấm thẻ là một số chẵn\rq\rq. Tính  xác suất để biến cố $C$ xảy ra (làm tròn đến hàng phần trăm).
	 \shortans{0,72}
	\loigiai{ 
		Dễ thấy $A$, $B$ là hai biến cố xung khắc và $\mathrm{P}(A)=\dfrac{\mathrm{C}_4^2}{\mathrm{C}_9^2}=\dfrac{1}{6}$, $\mathrm{P}(B)=\dfrac{\mathrm{C}_4^1\cdot  \mathrm{C}_5^1}{\mathrm{C}_9^{2}}=\dfrac{5}{9}$.\\
	Ta  có $\mathrm{P}(C)=\mathrm{P}(A\cup B)=\mathrm{P}(A)+\mathrm{P}(B)=\dfrac{1}{6}+\dfrac{5}{9}=\dfrac{13}{18}$.}
\end{ex}

% nguồn: Ok_TraLoiNgan_Lop11_Lan1
\begin{ex}%[1D9V2-4]%[ma:ce0f841b]
	Thanh có $4$ tấm thẻ được đánh số $1$, $3$, $4$, $7$. Thanh lấy ra $3$ trong $4$ thẻ và xếp chúng thành một hàng ngang một cách ngẫu nhiên để tạo thành một số có $3$  chữ số. Tính xác suất của biến cố $A$: \text{\lq\lq}Số tạo thành chia hết cho $2$ hoặc $3$\text{\rq\rq} (làm tròn kết quả đến hàng phần trăm).
	 \shortans{0,42}
	\loigiai{ 
	Số các số có $3$ chữ số có thể tạo thành là $24$ số.\\ 
	Gọi $B$ là biến cố \text{\lq\lq}Số tạo thành chia  hết cho $2$\text{\rq\rq}.\\
	Biến cố $B$ xảy ra khi chữ số  hàng đơn vị của số tạo thành là $4$.\\
	Có thể xếp được $3\cdot 2=6$ số chia hết  cho $2$.\\
	Do  đó $\mathrm{P}(B)=\dfrac{6}{24}$.\\
	Gọi $C$ là biến cố  \text{\lq\lq}Số tạo thành chia hết cho $3$\text{\rq\rq}.\\
	Biến cố $C$ xảy ra khi $3$ chữ số của số  tạo thành là $1;4;7$.\\
	Có thể xếp được $3\cdot 2=6$ số chia hết cho $3$.\\ 
	Do đó  $\mathrm{P}(C)=\dfrac{6}{24}$.\\
	Biến cố $BC$ xảy ra khi số tạo  thành chia hết cho $6$.\\
	Có $2$  kết quả thuận lợi cho biến cố $BC$. Do đó $\mathrm{P}(BC)=\dfrac{2}{24}$.\\
	Vậy  $\mathrm{P}(A)=\mathrm{P}(B\cup C)=\mathrm{P}(B)+\mathrm{P}(C)-\mathrm{P}(BC)=\dfrac{5}{12}\approx 0{,}42$.
	 }
\end{ex}

% nguồn: Ok_TraLoiNgan_Lop11_Lan1
\begin{ex}%[1D9C2-4]%[ma:d830d3ad]
	Hai vận động viên môn quần vợt A và B thi đấu trận chung kết một giải Grand Slam. Tay vợt nào thắng $3$ hiệp trước thi thắng trận. Trong mỗi hiệp đấu, xác suất để A thắng là $0{,}3$ và để B thắng là $0{,}7$. Tính xác suất để A thắng  trận chung kết đó (làm tròn đến hàng thứ $3$ sau dấu phẩy).
	 \shortans{0,16}
	\loigiai{ 
		Gọi $X$ là biến  cố A thấng trận đấu.\\
		Gọi $X_1$, $X_2$, $X_3$  tương ứng là biến cố A thắng B với tỉ số lần lượt là $3-0$, $3-1$, $3-2$.\\
		Suy ra  $X=X_1 \cup X_2 \cup X_3$ và $X_1$, $X_2$, $X_3$ đôi một xung khắc.\\
		Xét biến cố $X_1$: A thắng  B với tỉ số $3-0$.\\
		Khi đó  phải đấu $3$ hiệp và A thắng cả $3$ hiệp $\Rightarrow \mathrm{P}\left(X_1\right)=(0{,}3)^3=0{,}027$.\\
		Xét biến cố $X_2$: A  thắng B với tỉ số $3-1$.\\
		Khi đó phải đấu $4$ hiệp và B thắng duy nhất  $1$ hiệp trong $3$ hiệp đầu, A thắng $3$ hiệp còn lại\\
		 $\Rightarrow \mathrm{P}\left(X_2\right)=\mathrm{C}_3^1 \cdot 0{,}7 \cdot(0{,}3)^3=0{,}0567$.\\
		Xét biến cố $X_3$: A thắng  B với ti số $3-2$.\\
		Khi đó phải đấu $5$ hiệp và B thắng  $2$ trong $4$ hiệp đầu, A thắng $3$ hiệp còn lại
		$$\Rightarrow \mathrm{P}\left(X_3\right)=\mathrm{C}_4^2 \cdot(0{,}7)^2  \cdot(0,3)^3=0{,}07938.$$
		Vậy  xác suất để A thấng trận chung kết đó là
		 $$\mathrm{P}(X)=\mathrm{P}\left(X_1\right)+\mathrm{P}\left(X_2\right)+\mathrm{P}\left(X_3\right)=0{,}027+0{,}0567+0{,}07938 \approx 0{,}163.$$
	} 
\end{ex}
